% !TeX encoding = UTF-8
% P3CTeX optional CMP module integration test with default modules
\documentclass[diagrams,PAC,default]{P3CTeX}
\PECTeXconfig{
  cover=false,
  toc=false,
  no-plagi=false,
  data={
    student={Estudiant UOC},
    grade={Grau d'Enginyeria Inform\`atica},
    subj-fullname={Compiladors},
    subj-shortname={CMP},
    subj-code={PID\_00215089},
    ca-fullname={PAC 1: Visi\'o general i diagrames de Tombstone},
    ca-shortname={PAC1}
  }
}

\begin{document}

\section{Coexist\`encia amb els m\`oduls per defecte (\texttt{pxTAB}, \texttt{pxLST})}

\pxTable[style=header]{Fase, Sortida}{An\`alisi l\`exica, Testimonis}{An\`alisi sint\`actica, Arbre sint\`actic}

\section{Diagrames de Tombstone i Compiladors (\texttt{pxDIAG})}

\pxTombstoneCompile[title={Compilaci\'o en Java}]{Programa}{Java}{JVM}{M}
\pxTombstoneVM[title={M\`aquina virtual JVM}]{JVM}{Mx}
\pxTombstoneRun[title={Execuci\'o sobre la JVM}]{Programa}{Mx}[JVM]

\pxTombstoneBootstrap[title={Bootstrapping d'Oberon de 486 a P-III}]{Oberon}{486}{P-III}

\pxCompilerModel[title={Estructura interna d'un compilador}]

% A real figure environment with dynamic placement and an ordinary cross-reference.\n\pxTCompiler[float=htbp,caption={Floating compiler},label={diag:compiler-float}]{A}{B}{C}\nFloating diagram: \\ref{diag:compiler-float}.\n\end{document}
